% !TEX program = xelatex
\documentclass[11pt,a4paper]{article}
\usepackage{../../../00_common/preamble}

\title{2025年度 (AY2025) 同志社大学大学院 生命医科学研究科\\
医工学コース 入学試験問題「数学」解答例}
\author{}
\date{}

\begin{document}
\maketitle
\vspace{-3em}

%======================================================================
\section*{第1問}

\subsection*{前提整理}

$\sin x$ は実数全体で $C^\infty$ 級なので,$x=0$ のまわりで Taylor 展開できる.
$n$ 階導関数は
\[
(\sin x)^{(n)}=\sin\!\Bigl(x+\frac{n\pi}{2}\Bigr)
\]
であり,$x=0$ を代入すると
\[
(\sin x)^{(n)}\Big|_{x=0}=\sin\frac{n\pi}{2}
=\begin{cases}
0 & (n=2k:\ \text{偶数}),\\[2pt]
(-1)^{k} & (n=2k+1:\ \text{奇数}).
\end{cases}
\]
よって偶数次の係数はすべて $0$,奇数次 $n=2k+1$ の係数は
$\dfrac{(-1)^k}{(2k+1)!}$ となる.

\subsection*{展開}

\paragraph{Step 1: Maclaurin 級数を書き下す.}
Taylor の定理 $f(x)=\sum_{n=0}^{\infty}\frac{f^{(n)}(0)}{n!}x^{n}$ に上の係数を代入して
\[
\sin x=\sum_{k=0}^{\infty}\frac{(-1)^{k}}{(2k+1)!}\,x^{2k+1}
=x-\frac{x^{3}}{3!}+\frac{x^{5}}{5!}-\frac{x^{7}}{7!}+\cdots.
\]

\paragraph{Step 2: 収束範囲を確かめる.}
一般項を $a_k=\dfrac{x^{2k+1}}{(2k+1)!}$ とおくと,比は
\[
\left|\frac{a_{k+1}}{a_k}\right|
=\frac{|x|^{2}}{(2k+2)(2k+3)}\xrightarrow[k\to\infty]{}0
\]
であり,任意の $x\in\R$ で $0<1$ を満たす.ゆえに d'Alembert の判定法により
級数は絶対収束し,収束半径は $R=\infty$ である.剰余項についても
$|R_n(x)|\le\dfrac{|x|^{n+1}}{(n+1)!}\to0$ なので,展開は $-\infty<x<\infty$ の全域で
$\sin x$ に一致する.

\[
\ans{\
\displaystyle\sin x=\sum_{k=0}^{\infty}\frac{(-1)^{k}}{(2k+1)!}\,x^{2k+1}
=x-\frac{x^{3}}{3!}+\frac{x^{5}}{5!}-\frac{x^{7}}{7!}+\cdots
\quad(-\infty<x<\infty)\ }
\]

検算:$x=\frac{\pi}{6}$ で部分和 $x-\frac{x^3}{6}+\frac{x^5}{120}=0.4999996\ldots$ となり,
真値 $\sin\frac{\pi}{6}=0.5$ に一致する.$\checkmark$

%======================================================================
\section*{第2問}

\subsection*{前提整理}

積分領域 $D_n=\{(x,y):x^2+y^2\le n^2,\ x\ge0,\ y\ge0\}$ は,
原点中心・半径 $n$ の円板のうち第1象限にある四分円である.
$n\to\infty$ では第1象限全体 $\{x\ge0,\ y\ge0\}$ に広がる.
被積分関数 $e^{-x^2-y^2}>0$ なので,広義重積分として極座標で評価する.

\subsection*{計算}

\paragraph{Step 1: 極座標へ変換する.}
$x=r\cos\theta,\ y=r\sin\theta$ とおくと,$dxdy=r\,dr\,d\theta$,
領域 $D_n$ は $0\le r\le n,\ 0\le\theta\le\frac{\pi}{2}$ に対応する.
$x^2+y^2=r^2$ より
\[
\iint_{D_n}e^{-x^2-y^2}\,dxdy
=\int_{0}^{\pi/2}\!\!\int_{0}^{n}e^{-r^{2}}\,r\,dr\,d\theta.
\]

\paragraph{Step 2: 動径方向を積分する.}
$\dfrac{d}{dr}\Bigl(-\tfrac12 e^{-r^2}\Bigr)=r\,e^{-r^2}$ なので
\[
\int_{0}^{n}e^{-r^{2}}r\,dr
=\Bigl[-\tfrac12 e^{-r^{2}}\Bigr]_{0}^{n}
=\frac12\bigl(1-e^{-n^{2}}\bigr).
\]
角度方向は $\int_0^{\pi/2}d\theta=\dfrac{\pi}{2}$ だから
\[
\iint_{D_n}e^{-x^2-y^2}\,dxdy
=\frac{\pi}{2}\cdot\frac12\bigl(1-e^{-n^{2}}\bigr)
=\frac{\pi}{4}\bigl(1-e^{-n^{2}}\bigr).
\]

\paragraph{Step 3: 極限をとる.}
$n\to\infty$ で $e^{-n^2}\to0$ ゆえ
\[
\ans{\
\displaystyle\lim_{n\to\infty}\iint_{D_n}e^{-x^2-y^2}\,dxdy=\frac{\pi}{4}\ }
\]

\begin{rem}
この結果は Gauss 積分 $\int_0^\infty e^{-x^2}dx=\frac{\sqrt\pi}{2}$ の導出に等しい.
実際,第1象限全体での積分は
$\bigl(\int_0^\infty e^{-x^2}dx\bigr)^2=\bigl(\frac{\sqrt\pi}{2}\bigr)^2=\frac{\pi}{4}$
と直積に分離でき,上の値と一致する.
\end{rem}

検算:有限 $n=8$ について $\frac{\pi}{4}(1-e^{-64})=0.7853981634$ であり,
$\frac{\pi}{4}=0.7853981634$ と一致する.$\checkmark$

%======================================================================
\section*{第3問}

\subsection*{前提整理}

\[
A=\begin{pmatrix}1&2&1\\0&1&0\\2&-1&1\end{pmatrix}
\]
の逆行列を求める.まず正則性を行列式で確認し,掃き出し法で $A^{-1}$ を構成する.

\subsection*{計算}

\paragraph{Step 1: 行列式を計算する（正則性の確認）.}
第2行に沿って余因子展開すると($a_{21}=a_{23}=0$,$a_{22}=1$)
\[
\det A=1\cdot(-1)^{2+2}\det\begin{pmatrix}1&1\\2&1\end{pmatrix}
=1\cdot(1\cdot1-1\cdot2)=-1\neq0.
\]
よって $A$ は正則で,逆行列が存在する.

\paragraph{Step 2: 掃き出し法で $A^{-1}$ を求める.}
$[\,A\mid I\,]$ を行基本変形する.
\[
\left[\begin{array}{ccc|ccc}
1&2&1&1&0&0\\
0&1&0&0&1&0\\
2&-1&1&0&0&1
\end{array}\right]
\xrightarrow{R_3\to R_3-2R_1}
\left[\begin{array}{ccc|ccc}
1&2&1&1&0&0\\
0&1&0&0&1&0\\
0&-5&-1&-2&0&1
\end{array}\right]
\]
\[
\xrightarrow{R_3\to R_3+5R_2}
\left[\begin{array}{ccc|ccc}
1&2&1&1&0&0\\
0&1&0&0&1&0\\
0&0&-1&-2&5&1
\end{array}\right]
\xrightarrow{R_3\to -R_3}
\left[\begin{array}{ccc|ccc}
1&2&1&1&0&0\\
0&1&0&0&1&0\\
0&0&1&2&-5&-1
\end{array}\right]
\]
最後に第3列・第2列を上方向に消去する.
\[
\xrightarrow{R_1\to R_1-R_3}
\left[\begin{array}{ccc|ccc}
1&2&0&-1&5&1\\
0&1&0&0&1&0\\
0&0&1&2&-5&-1
\end{array}\right]
\xrightarrow{R_1\to R_1-2R_2}
\left[\begin{array}{ccc|ccc}
1&0&0&-1&3&1\\
0&1&0&0&1&0\\
0&0&1&2&-5&-1
\end{array}\right]
\]
左半が単位行列になったので,右半が逆行列である.
\[
\ans{\
A^{-1}=\begin{pmatrix}-1&3&1\\0&1&0\\2&-5&-1\end{pmatrix}\ }
\]

検算:
\[
A\,A^{-1}=\begin{pmatrix}1&2&1\\0&1&0\\2&-1&1\end{pmatrix}
\begin{pmatrix}-1&3&1\\0&1&0\\2&-5&-1\end{pmatrix}
=\begin{pmatrix}1&0&0\\0&1&0\\0&0&1\end{pmatrix}=I.
\]
たとえば $(1,1)$ 成分は $1\cdot(-1)+2\cdot0+1\cdot2=1$,
$(3,2)$ 成分は $2\cdot3+(-1)\cdot1+1\cdot(-5)=0$ となり,確かに単位行列.$\checkmark$

%======================================================================
\section*{第4問}

\subsection*{前提整理}

\[
B=\begin{pmatrix}4&1&0\\0&2&0\\0&0&3\end{pmatrix}
\]
が対角化可能かを調べ,可能ならば対角化する.方針は,固有値を求め,
各固有値の固有空間の次元(幾何的重複度)が代数的重複度に一致するかを確認することである.
$3$ 個の相異なる固有値が得られれば,その時点で対角化可能と結論できる.

\subsection*{計算}

\paragraph{Step 1: 固有値を求める.}
$B$ は上三角行列なので,固有方程式は対角成分から直ちに得られる:
\[
\det(B-\lambda I)
=(4-\lambda)(2-\lambda)(3-\lambda)=0
\quad\Longrightarrow\quad
\lambda=4,\ 2,\ 3.
\]
$3$ 個の固有値はすべて相異なる(各々が単根)ので,$B$ は\textbf{対角化可能}である.

\paragraph{Step 2: 各固有値の固有ベクトルを求める.}
$(B-\lambda I)\vv=\vzero$ を解く.

\smallskip
\noindent$\lambda=4$:
\[
B-4I=\begin{pmatrix}0&1&0\\0&-2&0\\0&0&-1\end{pmatrix},\quad
\begin{cases}v_2=0\\ v_3=0\end{cases}
\ \Rightarrow\ \vv_1=\begin{pmatrix}1\\0\\0\end{pmatrix}.
\]

\noindent$\lambda=2$:
\[
B-2I=\begin{pmatrix}2&1&0\\0&0&0\\0&0&1\end{pmatrix},\quad
\begin{cases}2v_1+v_2=0\\ v_3=0\end{cases}
\ \Rightarrow\ \vv_2=\begin{pmatrix}1\\-2\\0\end{pmatrix}.
\]

\noindent$\lambda=3$:
\[
B-3I=\begin{pmatrix}1&1&0\\0&-1&0\\0&0&0\end{pmatrix},\quad
\begin{cases}v_1+v_2=0\\ -v_2=0\end{cases}
\ \Rightarrow\ v_1=v_2=0,\ \ \vv_3=\begin{pmatrix}0\\0\\1\end{pmatrix}.
\]

\paragraph{Step 3: 変換行列を組み立てる.}
固有ベクトルを列に並べた $P=[\,\vv_1\ \vv_2\ \vv_3\,]$ により $P^{-1}BP$ は
対応する固有値を並べた対角行列になる.
\[
\ans{\
P=\begin{pmatrix}1&1&0\\0&-2&0\\0&0&1\end{pmatrix},\qquad
P^{-1}BP=\begin{pmatrix}4&0&0\\0&2&0\\0&0&3\end{pmatrix}\ }
\]

検算:$\det P=-2\neq0$ より $P$ は正則.
$BP$ を計算すると
\[
BP=\begin{pmatrix}4&1&0\\0&2&0\\0&0&3\end{pmatrix}
\begin{pmatrix}1&1&0\\0&-2&0\\0&0&1\end{pmatrix}
=\begin{pmatrix}4&2&0\\0&-4&0\\0&0&3\end{pmatrix}
=\begin{pmatrix}1&1&0\\0&-2&0\\0&0&1\end{pmatrix}
\begin{pmatrix}4&0&0\\0&2&0\\0&0&3\end{pmatrix}=PD,
\]
すなわち $BP=PD$,したがって $P^{-1}BP=D=\mathrm{diag}(4,2,3)$ が成り立つ.$\checkmark$

\vspace{1.5em}
\begin{rem}
固有ベクトルの取り方には定数倍・符号の自由度があるため,$P$ の各列を
スカラー倍したものも正解である(対応する $D$ は変わらない).
列の並べ替えを行えば $D$ の対角成分の順序もそれに合わせて変わる.
\end{rem}

\end{document}
